% De Haruspicum Responsis (de-haruspicum-responsis) --- English translation
% Translated by Claude (Anthropic) from the Latin (Perseus).
% Style guide: STYLE.md.

\ciceroSection{1} Yesterday, senators, when both your dignity and the great press of Roman knights present (to whom the Senate was being granted access) had stirred me deeply, I judged that I had to put down P. Clodius's shameless impudence: when he was hindering the cause of the publicans by the most foolish questionings, doing service for P. Tullio the Syrian, and selling himself --- though he had been bought up entirely --- even before your eyes. So I held in check that frenzied and exulting man as soon as I threatened him with the danger of a trial: with two opening words I crushed all the gladiator's onset and ferocity.

\ciceroSection{2} And yet, ignorant who the consuls were, ashen and hot-flushed, he flung himself out of the Curia all at once, with certain broken and now empty threats, and with those terrors of the Pisonian and Gabinian period: whom, when he had gone out, I began to follow, and I had the greatest reward both from the rising-up of every one of you and from the company of the publicans. But suddenly out of his wits, without his own face, without colour, without voice, he stood still; then he looked back, and as soon as he caught sight of Cn. Lentulus the consul, he collapsed almost at the threshold of the Curia --- by recollection, I suppose, of his Gabinius, and by longing for Piso. Of his unbridled, headlong frenzy what shall I say? Or can he be more grievously wounded by my words than he was, in the very act, dispatched and butchered by that gravest of men, P. Servilius? Whose force and signal, almost divine gravity, even if I could now match, still I do not doubt that the weapons hurled by an enemy will seem lighter and duller than those his father's colleague let fly.

\ciceroSection{3} Yet I wish to set out the reasoning of my act to those who, yesterday, supposed me to have been carried by pain and to have gone almost further in anger than the considered reason of a wise man would have demanded. I did nothing in anger, nothing with an unmastered mind, nothing that was not long thought through and meditated long before. For I, senators, have always declared that I am a foe to two men: who, when they ought to have defended me and the commonwealth, when they could have saved us, when they were called to the consular duty by the very ensigns of that command, and to my safety not only by your authority but by your prayers, first deserted us, then betrayed us, lastly attacked us; and by the rewards of an unspeakable bargain wished me, together with the commonwealth, to be utterly crushed and snuffed out. Those punishments which they could not, under their own bloody and fatal command, either keep from the walls of allies or carry into the cities of enemies --- the cutting-down, the burning, the overturning, the depopulation, the wasteland --- these, with their plunder, they brought into all my houses and fields.

\ciceroSection{4} On these Furies and firebrands, on these, I say, ruinous prodigies and almost pestilences of this empire, I declare that I have undertaken an inexpiable war: not as great as my own pain and that of mine demanded, but as great as your pain and the pain of all good men demanded. Toward Clodius, indeed, my hatred is no greater today than it was on the day when I learned that he had been singed by the most religious fires, sent away in woman's dress out of an unholy seduction and out of the house of the pontifex maximus. Then, I say, then I saw, and long before I had foreseen, what storm was being raised, what tempest was hanging over the commonwealth. I saw that wickedness, so monstrous, the audacity so savage of a youth raving and noble and wounded, could not be confined within the bounds of peace: that evil would burst out at some time, if it had gone unpunished, to the destruction of the city. Not much, in fact, was added afterwards to my hatred.

\ciceroSection{5} For he did nothing against me out of hatred of me, but out of hatred of severity, hatred of standing, hatred of the commonwealth: he no more violated me than he violated the Senate, the Roman knights, all good men, all Italy. Nor was he in the end more wickedly disposed against me than against the immortal gods themselves. For he violated them by a wickedness no man before him had committed; toward me he was of the same temper which his bosom-friend Catiline, had he won, would have shown. So I never thought he should be prosecuted by me --- no more than that block of wood, of which we should not know what kind of man it was, did he not himself say he was a Ligurian. Why should I pursue this animal and beast, fattened on the fodder and acorns of my enemies? Who, if he sensed by what wickedness he has bound himself, I do not doubt is the most miserable of men; but if he does not see it, there is danger that he will defend himself by the excuse of stupor.

\ciceroSection{6} It is added too that, by the expectation of all, this victim seems to have been vowed and assigned to that bravest and most distinguished man, T. Annius. To snatch from him a glory already pledged and destined --- when I myself by his work have recovered both my standing and my safety --- would be very unjust. For just as that famous P. Scipio seems to me to have been born for the destruction and ruin of Carthage, who alone overturned by a kind of fated coming that city which had been by many commanders besieged, attacked, shaken, almost taken, so T. Annius seems to have been born for the suppressing, snuffing out, and razing to the ground of this pestilence, and granted as it were by a divine gift to the commonwealth. He alone has understood by what method an armed citizen --- one who with stones, with iron, drives some men to flight, holds others shut up at home; who terrifies the whole city, the Curia, the Forum, all the temples, with slaughter and burnings --- ought not only to be conquered but bound.

\ciceroSection{7} From a man so great, and one who has so deserved of me and of his country, I will never voluntarily snatch this defendant in particular --- whose enmities he has not only undertaken on my behalf but has even sought out. But if even now, ensnared as he is in the dangers of every law, entangled in the hatred of all good men, caught in the now not distant expectation of punishment, he is yet borne forward fumbling and tries, blocked though he is, to make some onset against me, I will resist; and either with Milo's consent or with his help I will repel his attempt: just as yesterday, when standing he silently threatened me, I touched only with my voice the beginning of the law and the trial. He sat down: I fell silent. Had he named the day, as he had hinted, I would have brought it about that the third day from the praetor was named for him at once. And let him so govern himself, and consider this: if he is content with the wickednesses he has committed, he is now consecrated to Milo; if he aims any weapon at me, I shall at once snatch up the arms of trials and laws.

\ciceroSection{8} A little while ago, senators, he held a contio which has been reported to me in full. The whole substance and drift of that contio, hear first; then, when you have laughed at the man's shamelessness, you shall hear from me about the whole. Of religion and the rites and ceremonies Clodius spoke in the contio, senators: Publius, I say, Publius Clodius complained that the rites and religious observances were being neglected, violated, polluted! No wonder it seems ridiculous to you: even his own contio laughed at the man, who, as he himself is wont to boast, has been transfixed by two hundred decrees of the Senate, all passed against him on behalf of religion --- and that man, who has brought defilement onto the cushions of the Good Goddess, who has violated those rites which it is forbidden for the eyes of a man to look on even unaware, not only with a man's gaze but with debauchery and seduction, complains in the contio that religion has been neglected.

\ciceroSection{9} And so now his next contio is awaited on chastity. For what difference does it make whether, having been driven from the most religious altars, he should complain of religion --- or whether, just emerged from his sister's bedroom, he should defend modesty and chastity? He read out in the contio this fresh response of the haruspices about the rumbling, in which, among many other things, this also is written, as you have heard, that ``sacred and religious places are being held profane.'' And he said in this matter that my house had been consecrated by that most religious priest, P. Clodius.

\ciceroSection{10} I am glad that out of this whole portent, which perhaps has been the most grave to be reported to this order in many years, a cause has been given to me of speaking, not only just but even necessary; for you will find that out of this whole prodigy and response we are forewarned of his wickedness and frenzy and of the impending greatest dangers almost by the very voice of Jupiter Best and Greatest.

\ciceroSection{11} But first I will expiate the religion of my house, if I shall be able to do so truly and without doubt on anyone's part; if even the slightest scruple shall seem to remain to anyone, not only with patience but with a willing mind I will obey the portents of the immortal gods and religion. But what house in this city is so empty of and so pure from any such suspicion of religion? Although your houses, senators, and those of the rest of the citizens are for the most part by far free from religion, yet my one house has been freed by all judgements: the only one in this city. For I appeal to you, Lentulus, and to you, Philippus. From this haruspical response the Senate decreed that you should refer to this order on sacred and religious places. Can you bring up my house, which, as I have said, alone in this whole city has been freed by every religion and by every judgement? Whose first enemy himself, in that storm and night of the commonwealth, when he had inscribed the rest of his crimes by that filthy stylus dipped in the mouth of Sex. Clodius, did not touch a single letter of religion; then the same house the Roman people, whose is the highest power in all matters, by the centuriate assembly, by the votes of every age and order, ordered to be of the same right as it had been; afterwards you, senators --- not because the matter was doubtful but so that to this Fury, if he should remain longer in the city he so longed to destroy, the very voice of speech might be denied --- decreed that the religion of my house be referred to the college of the pontiffs.

\ciceroSection{12} What religion is so great that we are not freed from it, in our doubts and in the gravest scruples, by the response and word of a single P. Servilius or M. Lucullus? Concerning public sacrifices, the greatest games, the rites of the household gods and Mother Vesta, that very sacrifice which is performed for the safety of the Roman people, which since the founding of Rome was first violated by the wickedness of this single chaste guardian of religion --- whatever three pontiffs had decided has always seemed to the Roman people, to the Senate, to the immortal gods themselves, holy enough, august enough, religious enough. But on my house, P. Lentulus, consul and pontiff; P. Servilius, M. Lucullus, Q. Metellus, M'. Glabrio, M. Messalla, L. Lentulus, flamen Martialis; P. Galba, Q. Metellus Scipio, C. Fannius, M. Lepidus, L. Claudius rex sacrorum, M. Scaurus, M. Crassus, C. Curio, Sex. Caesar flamen Quirinalis, Q. Cornelius, P. Albinovanus, Q. Terentius, the lesser pontiffs --- having heard the cause, with it pleaded twice, with the greatest assembly of the most distinguished and wisest citizens standing by --- with one mind all freed it from every religious obligation.

\ciceroSection{13} I deny that, since the establishment of the rites --- whose antiquity is the same as the city's --- in any matter whatever, not even concerning the chastity of the Vestal virgins, has so full a college given judgement. Although for the discussion of a crime it matters that as many as possible be present (for that is the rule of pontifical interpretation, since the same men hold the power of judges), the explanation of religion can be rightly done even by one experienced pontiff (a thing which in a capital trial is harsh and unjust); yet you will find this: that more pontiffs have judged on my house than ever judged on the rites of the virgins. The day after, the Senate, in fullest attendance, with you, Lentulus, the consul-elect, leader of the opinion, on the motion of P. Lentulus and Q. Metellus the consuls, when all the pontiffs of this order were present and when others who excelled in honours of the Roman people had spoken many words about the college's verdict, and when all of the same were present at the drawing-up, decreed that my house was seen to be freed from religion by the judgement of the pontiffs.

\ciceroSection{14} On this place, then, sacred above all, do the haruspices seem to be speaking, on the place which alone of all private places has this peculiar right, that by the very men who preside over the rites it has been judged not sacred? But bring up the matter, as you must do under the senatorial decree. Either the inquiry will be granted to you, who first delivered an opinion on this house and freed it from every religious obligation; or the Senate itself will judge, which, with that one priest of the rites alone dissenting, in fullest attendance previously judged; or, what surely will happen, it will be referred to the pontiffs, to whose authority, faith, and prudence our forefathers entrusted both private and public rites and religious matters. What then can they judge other than they have judged? Many are the houses in this city, senators, and I dare say almost all by the best of right, but yet by private right, by hereditary right, by right of authority, by right of mancipium, by right of bond: I deny that any other house exists fortified by the same private right as the one of best title, and at the same time by every public, human, and divine right outstanding;

\ciceroSection{15} since first it is built on the authority of the Senate at public expense, and then it has been fortified and hedged around by so many decrees of the Senate against the wicked violence of this gladiator. The first commission was given last year to those magistrates --- to whom in the gravest dangers the whole commonwealth is wont to be entrusted --- to see that I might be allowed to build without violence; then, when he had brought ruin to my dwelling with stones and fire and steel, the Senate decreed that those who had done it should be held under the law on violence, which is against those who attack the whole commonwealth. And on your motion --- O bravest and best consuls within memory! --- the same Senate, in fullest attendance, decreed that whoever should have violated my house would be acting against the commonwealth.

\ciceroSection{16} I deny that there exist on any public work, on any monument, on any temple, as many decrees of the Senate as on my house, which the Senate has held to be the only one since the city's founding to be built from the treasury, to be freed by the pontiffs, to be defended by the magistrates, to be punished by the courts. To P. Valerius for his great services to the commonwealth a public house was given on the Velia; to me, on the Palatine, restored. To him, the place; to me, even the walls and the roof. To him, what he himself should defend by private right; to me, what every magistrate should defend publicly. Which things, indeed, if I held either of myself or by others, I would not proclaim before you, lest I should seem to glory too much; but since they are given me by you, since they are being attempted by the tongue of him by whose hand they were before overthrown, you who restored me and my children with your own hands --- I am not speaking of my own deeds but of yours; nor do I fear that this proclamation of mine, of your kindnesses, may seem ungrateful rather than arrogant.

\ciceroSection{17} And yet if, after I have undergone such labours for the common safety, in refuting the slanders of wicked men a certain pang of spirit should at any time carry me away to glory, who would not pardon me? For yesterday I saw a certain man muttering, who, they said, was saying that I could not be borne, because, when I was asked by this same most foul parricide what city I was a citizen of, I answered, with you and the Roman knights approving, that I was the citizen of that city which had not been able to do without me. He, I think, groaned. What then was I to answer? --- I ask of him himself who cannot bear me. That I am a Roman citizen? It would have been a literal answer. Or to be silent? It would have been to abandon the matter. Can any man, engaged in great affairs amid envy, answer the insults of an enemy gravely enough without his own praise? But he himself not only answers whatever he can when challenged, but actually rejoices that he is reminded by his friends what he should answer.

\ciceroSection{18} But since my own case is dispatched, let us see now what the haruspices say. For I confess that I have been deeply stirred both by the magnitude of the portent and by the gravity of the response and by the one and consistent voice of the haruspices; nor am I one who, if perhaps I seem to be more than the rest who are equally busy with me engaged in the study of letters, take pleasure in or use at all those letters which would deter our minds and call them away from religion. I have, indeed, our forefathers as authors and masters of the cultivating of religion, whose wisdom seems to me to have been so great that those are wise enough and beyond who can not, I will not say match their wisdom, but at least see how great it was: who held that fixed and solemn rites are contained in the pontificate, the auspices for things well done in augury, the old prophecies of fate in the books of the Apolline seers, and the expiations of portents in the discipline of the Etruscans. Which is so great that within our memory it foretold to us, not obscurely but a little while before, first the dire beginnings of the Italic war; afterwards the almost-extreme crisis of the Sullan and Cinnan period; and then this recent conspiracy to set the city on fire and to destroy the empire.

\ciceroSection{19} Then, if I had any leisure, I have learned besides that many learned and wise men have both said and left in writing many things on the godhead of the immortal gods; which, although I see they are written down divinely, are yet of such a kind that our forefathers seem to have taught these men, not to have learned from them. For who is so out of his wits that, when he has looked up to the sky, does not feel that there are gods, and that those things are formed by such a great mind that scarcely anyone by any art could trace through the order and the necessary connection of things --- thinks all this happens by chance? Or, when he has understood that there are gods, does not understand that this great empire was born and increased and held by their godhead? However much we may, senators, love ourselves, yet not in number do we surpass the Spaniards, nor in strength the Gauls, nor in cleverness the Phoenicians, nor in arts the Greeks, nor finally in that very domestic and inborn sense of this nation and this land the Italians and Latins themselves --- but by piety and religion and by this one wisdom, that we have seen all things to be ruled and steered by the godhead of the gods, we have surpassed all peoples and nations.

\ciceroSection{20} So, that I may not say more on a matter least doubtful, lend your attention, and bring not only your ears but your minds to the haruspices' voice: ``\textit{Because in the Latinian field a noise was heard with a rumbling.}'' I leave aside the haruspices, I leave aside that ancient discipline, handed down to Etruria, as men say, by the immortal gods themselves: cannot we ourselves be haruspices? In the field nearby and suburban there was heard a hidden noise and horrible rumbling of arms. Who is there of those Giants, whom the poets tell us made war on the immortal gods, so impious that he would not confess that by such a new and great motion the gods are foretelling and proclaiming something great to the Roman people? Of the matter it is written: ``\textit{There are postilia owing to Jupiter, Saturn, Neptune, Tellus, the heavenly gods.}''

